\lecture{5}
\title{Inverse Kinematics via Optimization}

\begin{document}

\subsection{Inversion as Optimization: Constraints}

Recall from \href{/6.4210/lectures/4}{Lecture 4}
that our approach to solving ``inverse'' problems
is to take the (damped, possibly) pseudo-inverse,
because it nicely satisfies the following properties.

\begin{theorem}{Pseudo-Inverse Approximation}
    For any matrix $J$ and vector $b$ (with appropriate dimensions),
    the vector $x^* := J^+b$ is the vector with minimum magnitude
    across all vectors that minimize $\|Jx - b\|$.  In particular,
    \begin{itemize}
        \item If the column space of $J$ does not contain $b$,
        then $x^* := J^+b$ is such that $\|Jx - b\|$ is minimized.
        \item If the column space of $J$ contains $b$
        (that is, $\|Jx - b\|$ can be zero),
        then $x^* := J^+b$ is the solution to $Jx = b$ minimizing $\|x\|$.
    \end{itemize}
\end{theorem}

We can go even further, though, by phrasing this as an optimization problem.

\begin{quote}
    \textbf{Quadratic Program.}
    Find the vector $x^*$ with minimum magnitude
    that minimizes $\|Jx^* - b\|^2$.
\end{quote}

Viewing this as a \textbf{quadratic program}
(see \href{/6.1220/lectures/11}{6.1220 Notes} for context),
we can also consider adding \textbf{constraints} to this problem.
For example, we might want to specify that
$|\mathrm{d}q_1| \leq 10$, $|\mathrm{d}q_2| \leq 10$ and
$\mathrm{d}q_1^2 + \mathrm{d}q_2^2 \leq 150$ to prevent
exploding velocities.

Notice that this is \emph{not} just as simple as
taking the output of the pseudo-inverse
and clipping it to a constrained region!
\begin{center}
    \frame{\includegraphics[width=0.35\textwidth]{lecture-05/clipping.png}}
\end{center}
As the image above indicates,
the solution achieved by clipping (shown in red)
can oftentimes be less optimal
than the solution achieved by optimizing the constrained quadratic program
(shown in blue).

\subsection{Inversion as Optimization: Joint Centering}

Now suppose that our system is highly redundant---there are
infinitely many $\mathrm{d}q$ that achieve
$J(q) \mathrm{d}q = \mathrm{d}X^B$.
Can we get something out of the fact that
our null space is nontrivial?

Well, the pseudo-inverse exploits this
by specifically aiming for the solution $\mathrm{d}q$ with smallest magnitude.

But maybe, more smartly, we can specify a \textbf{preferred position} $q^d$.
Then if our current state is $q$, we might imagine a spring of stiffness $K$
that induces a desired velocity of $K(q^d - q)$.

Since the null space is nontrivial,
we can \emph{project} this desired velocity $K(q^d - q)$ into the null space.
This lets us freely add the projection into $\mathrm{d}q$
without affecting the fact that $J(q) \mathrm{d}q = \mathrm{d}X^B$.
See the (warning: AI-generated) widget below.

\includewidget{lecture-05/joint-centering}

This is fairly widely used, especially for robots that are highly redundant.

\subsection{Trajectory Optimization}

A common constraint that one might use
is that the robot must \emph{avoid collisions}.
In the very simple example of a 2D simplified two-rod Kuka IIWA,
here's how that constraint looks in the configuration space $(q_1, q_2)$.
\begin{center}
    \frame{\includegraphics[width=0.5\textwidth]{lecture-05/no-collision.png}}
\end{center}
Notably, this constraint is \textbf{non-convex},
so it cannot simply be solved by traditional quadratic program solvers.
(For example, gradient descent can get stuck in a local minimum.)

One might define a \textbf{trajectory optimization} problem like so:
\begin{quote}
    \textbf{Trajectory Optimization.}
    Minimize $\sum_{n = 0}^{N - 1} \| q[n + 1] - q[n] \|^2$
    such that $q[0] = q_{\mathrm{start}}$, $q[N] = q_{\mathrm{goal}}$,
    and the path $q[0] \to q[1] \to \dots \to q[N]$
    does not collide with anything.
\end{quote}
An optimizer might then start with an initial guess
and gradually improve the guess, in hopes that
the local minimum that it converges towards
does actually satisfy the nonconvex constraints.
\begin{center}
    \frame{\includegraphics[width=0.5\textwidth]{lecture-05/trajectory-optimization.png}}
\end{center}
Given a set of points $\{q[n]\}_{n = 0}^N$,
we can define a path via a \textbf{B-spline}, for example.
\begin{center}
    \frame{\includegraphics[width=0.35\textwidth]{lecture-05/b-spline.png}}
\end{center}
One nice property of B-splines is that the curve
\emph{always stays within the convex hull of the control points}.
But one bad property of B-splines is that
even if the control points satisfy the constraints,
the path itself might not.
\begin{center}
    \frame{\includegraphics[width=0.25\textwidth]{lecture-05/failure-mode.png}}
\end{center}
Next lecture, we'll discuss alternative (non-optimization-based) approaches
to motion planning.

\end{document}